% GENERATED FILE. Edit canonical lessons, then run npm run generate:books.
% canonical-source-hash: fnv1a64:aeb7e41c885c462f
% canonical-lessons: ML-C511-pravacanam, ML-C511-jagrata, ML-C511-usnataramgam, ML-C511-culalikkarru, ML-C511-pottiviluka

\chapter{Forecasts and Storms}
\label{ch:ml-forecasts-and-storms}
\hlchaptermodality{511}

\begin{chapteropening}
\textbf{By the end of this chapter:} \emph{I can say a forecast, vigilance, a heatwave, a cyclone and to snap and fall.}

Complete the last lesson of chapter 511: I can say a forecast, vigilance, a heatwave, a cyclone and to snap and fall.
\end{chapteropening}

\section[\texorpdfstring{\ml{പ്രവചനം} (pravacanaṁ)}{പ്രവചനം (pravacanaṁ)}]{\ml{പ്രവചനം} (pravacanaṁ) — a forecast; a prediction}

\label{lesson:ML-C511-pravacanam}



\begin{quote}
Before the new one: say the Malayalam for to grow cold, then the Malayalam for to rumble.
\end{quote}

\subsection*{You'll want to know: \ml{പ്രവചനം}}
\textbf{\ml{പ്രവചനം}} — \emph{pravacanaṁ} — \textquotedblleft{}a forecast; a prediction\textquotedblright{}.

\begin{quote}
\ml{നാളെ} \ml{കനത്ത} \ml{മഴ} \ml{ഉണ്ടാകുമെന്നാണ്} \ml{പ്രവചനം}.
\end{quote}

\begin{quote}
\emph{nāḷe kanatta maḻa uṇṭākumennāṇŭ pravacanaṁ.}
\end{quote}

\begin{quote}
The forecast is that there will be heavy rain tomorrow.
\end{quote}

A weather forecast is a \ml{കാലാവസ്ഥാ} \ml{പ്രവചനം}.

\begin{cousinweb}
\textbf{\ml{പ്രവചനം}} (\emph{pravacanaṁ}) is borrowed from Sanskrit \emph{pravacana}, \textquotedblleft{}speaking out, announcement\textquotedblright{} (Monier-Williams).
\end{cousinweb}

\begin{grammarlens}[title={What you've built}]
One more word for getting around.
\end{grammarlens}

\subsection*{Your turn}
\begin{itemize}
\raggedright
  \item \emph{Say it:} \emph{pravacanaṁ}
  \item \emph{Say it:} \emph{pravacanaṁ}, once more
  \item \emph{Say it:} \emph{nāḷe kanatta maḻa uṇṭākumennāṇŭ pravacanaṁ.} — The forecast is that there will be heavy rain tomorrow.
  \item \emph{Recall:} say the Malayalam for to grow cold, then the Malayalam for to rumble, then say \emph{pravacanaṁ} again
\end{itemize}

\begin{tcolorbox}[breakable,colback=teal!4,colframe=teal!35!black,title={Before you move on}]
What does \ml{പ്രവചനം} mean? (\textquotedblleft{}A forecast; a prediction\textquotedblright{}.) Say it once more.
\end{tcolorbox}
\section[\texorpdfstring{\ml{ജാഗ്രത} (jāgrata)}{ജാഗ്രത (jāgrata)}]{\ml{ജാഗ്രത} (jāgrata) — vigilance; an alert}

\label{lesson:ML-C511-jagrata}



\begin{quote}
Before the new one: say the Malayalam for to rumble, then the Malayalam for a forecast.
\end{quote}

\subsection*{You'll want to know: \ml{ജാഗ്രത}}
\textbf{\ml{ജാഗ്രത}} — \emph{jāgrata} — \textquotedblleft{}vigilance; an alert\textquotedblright{}.

\begin{quote}
\ml{മഴ} \ml{കാരണം} \ml{തീരദേശത്ത്} \ml{ജാഗ്രത} \ml{വേണം}.
\end{quote}

\begin{quote}
\emph{maḻa kāraṇaṁ tīradēśattŭ jāgrata vēṇaṁ.}
\end{quote}

\begin{quote}
Because of the rain, the coast must stay alert.
\end{quote}

Weather offices give a \ml{ജാഗ്രതാ} \ml{നിർദ്ദേശം}, an alert notice.

\begin{cousinweb}
\textbf{\ml{ജാഗ്രത}} (\emph{jāgrata}) is borrowed from Sanskrit \emph{jāgrat}, \textquotedblleft{}being awake\textquotedblright{} (Monier-Williams).
\end{cousinweb}

\begin{grammarlens}[title={What you've built}]
One more word for getting around.
\end{grammarlens}

\subsection*{Your turn}
\begin{itemize}
\raggedright
  \item \emph{Say it:} \emph{jāgrata}
  \item \emph{Say it:} \emph{jāgrata}, once more
  \item \emph{Say it:} \emph{maḻa kāraṇaṁ tīradēśattŭ jāgrata vēṇaṁ.} — Because of the rain, the coast must stay alert.
  \item \emph{Recall:} say the Malayalam for to rumble, then the Malayalam for a forecast, then say \emph{jāgrata} again
\end{itemize}

\begin{tcolorbox}[breakable,colback=teal!4,colframe=teal!35!black,title={Before you move on}]
What does \ml{ജാഗ്രത} mean? (\textquotedblleft{}Vigilance; an alert\textquotedblright{}.) Say it once more.
\end{tcolorbox}
\section[\texorpdfstring{\ml{ഉഷ്ണതരംഗം} (uṣṇataraṁgaṁ)}{ഉഷ്ണതരംഗം (uṣṇataraṁgaṁ)}]{\ml{ഉഷ്ണതരംഗം} (uṣṇataraṁgaṁ) — a heatwave}

\label{lesson:ML-C511-usnataramgam}



\begin{quote}
Before the new one: say the Malayalam for a forecast, then the Malayalam for vigilance.
\end{quote}

\subsection*{You'll want to know: \ml{ഉഷ്ണതരംഗം}}
\textbf{\ml{ഉഷ്ണതരംഗം}} — \emph{uṣṇataraṁgaṁ} — \textquotedblleft{}a heatwave\textquotedblright{}.

\begin{quote}
\ml{ഉഷ്ണതരംഗം} \ml{കാരണം} \ml{പള്ളിക്കൂടങ്ങൾ} \ml{അടച്ചു}.
\end{quote}

\begin{quote}
\emph{uṣṇataraṁgaṁ kāraṇaṁ paḷḷikkūṭaṅṅaḷ aṭaccu.}
\end{quote}

\begin{quote}
Because of the heatwave the schools were closed.
\end{quote}

\ml{ഉഷ്ണം}, heat, plus \ml{തരംഗം}, a wave.

\begin{cousinweb}
\textbf{\ml{ഉഷ്ണതരംഗം}} (\emph{uṣṇataraṁgaṁ}) is a modern Malayalam compound built from two Sanskrit words, \emph{uṣṇa}, \textquotedblleft{}hot\textquotedblright{}, and \emph{taraṅga}, \textquotedblleft{}a wave\textquotedblright{} (Monier-Williams).
\end{cousinweb}

\begin{grammarlens}[title={What you've built}]
One more word for getting around.
\end{grammarlens}

\subsection*{Your turn}
\begin{itemize}
\raggedright
  \item \emph{Say it:} \emph{uṣṇataraṁgaṁ}
  \item \emph{Say it:} \emph{uṣṇataraṁgaṁ}, once more
  \item \emph{Say it:} \emph{uṣṇataraṁgaṁ kāraṇaṁ paḷḷikkūṭaṅṅaḷ aṭaccu.} — Because of the heatwave the schools were closed.
  \item \emph{Recall:} say the Malayalam for a forecast, then the Malayalam for vigilance, then say \emph{uṣṇataraṁgaṁ} again
\end{itemize}

\begin{tcolorbox}[breakable,colback=teal!4,colframe=teal!35!black,title={Before you move on}]
What does \ml{ഉഷ്ണതരംഗം} mean? (\textquotedblleft{}A heatwave\textquotedblright{}.) Say it once more.
\end{tcolorbox}
\section[\texorpdfstring{\ml{ചുഴലിക്കാറ്റ്} (cuḻalikkāṟṟŭ)}{ചുഴലിക്കാറ്റ് (cuḻalikkāṟṟŭ)}]{\ml{ചുഴലിക്കാറ്റ്} (cuḻalikkāṟṟŭ) — a cyclone}

\label{lesson:ML-C511-culalikkarru}



\begin{quote}
Before the new one: say the Malayalam for vigilance, then the Malayalam for a heatwave.
\end{quote}

\subsection*{You'll want to know: \ml{ചുഴലിക്കാറ്റ്}}
\textbf{\ml{ചുഴലിക്കാറ്റ്}} — \emph{cuḻalikkāṟṟŭ} — \textquotedblleft{}a cyclone\textquotedblright{}.

\begin{quote}
\ml{ചുഴലിക്കാറ്റിൽ} \ml{പല} \ml{മരങ്ങളും} \ml{വീണു}.
\end{quote}

\begin{quote}
\emph{cuḻalikkāṟṟil pala maraṅṅaḷuṁ vīṇu.}
\end{quote}

\begin{quote}
Many trees fell in the cyclone.
\end{quote}

\ml{ചുഴലി}, whirling, plus \ml{കാറ്റ്}, wind.

\begin{cousinweb}
\textbf{\ml{ചുഴലിക്കാറ്റ്}} (\emph{cuḻalikkāṟṟŭ}) is a native Dravidian compound of \ml{ചുഴലി}, whirling (DEDR 2698), and \ml{കാറ്റ്}, wind (DEDR 1481).
\end{cousinweb}

\begin{grammarlens}[title={What you've built}]
One more word for getting around.
\end{grammarlens}

\subsection*{Your turn}
\begin{itemize}
\raggedright
  \item \emph{Say it:} \emph{cuḻalikkāṟṟŭ}
  \item \emph{Say it:} \emph{cuḻalikkāṟṟŭ}, once more
  \item \emph{Say it:} \emph{cuḻalikkāṟṟil pala maraṅṅaḷuṁ vīṇu.} — Many trees fell in the cyclone.
  \item \emph{Recall:} say the Malayalam for vigilance, then the Malayalam for a heatwave, then say \emph{cuḻalikkāṟṟŭ} again
\end{itemize}

\begin{tcolorbox}[breakable,colback=teal!4,colframe=teal!35!black,title={Before you move on}]
What does \ml{ചുഴലിക്കാറ്റ്} mean? (\textquotedblleft{}A cyclone\textquotedblright{}.) Say it once more.
\end{tcolorbox}
\section[\texorpdfstring{\ml{പൊട്ടിവീഴുക} (poṭṭivīḻuka)}{പൊട്ടിവീഴുക (poṭṭivīḻuka)}]{\ml{പൊട്ടിവീഴുക} (poṭṭivīḻuka) — to snap and fall}

\label{lesson:ML-C511-pottiviluka}



\begin{quote}
Before the new one: say the Malayalam for a heatwave, then the Malayalam for a cyclone.
\end{quote}

\subsection*{You'll want to know: \ml{പൊട്ടിവീഴുക}}
\textbf{\ml{പൊട്ടിവീഴുക}} — \emph{poṭṭivīḻuka} — \textquotedblleft{}to snap and fall\textquotedblright{}.

\begin{quote}
\ml{കാറ്റിൽ} \ml{മരക്കൊമ്പ്} \ml{പൊട്ടിവീണു}.
\end{quote}

\begin{quote}
\emph{kāṟṟil marakkompŭ poṭṭivīṇu.}
\end{quote}

\begin{quote}
In the wind a branch snapped and fell.
\end{quote}

Past: \ml{പൊട്ടിവീണു}, \emph{poṭṭivīṇu}.

\ml{പൊട്ടി}, bursting or snapping, plus \ml{വീഴുക}, to fall.

\begin{cousinweb}
\textbf{\ml{പൊട്ടിവീഴുക}} (\emph{poṭṭivīḻuka}) is a native Dravidian compound of \ml{പൊട്ടുക}, to snap or burst (DEDR 4490), and \ml{വീഴുക}, to fall (DEDR 5430).
\end{cousinweb}

\begin{grammarlens}[title={What you've built}]
One more word for getting around.
\end{grammarlens}

\subsection*{Your turn}
\begin{itemize}
\raggedright
  \item \emph{Say it:} \emph{poṭṭivīḻuka}
  \item \emph{Say it:} \emph{poṭṭivīḻuka}, once more
  \item \emph{Say it:} \emph{kāṟṟil marakkompŭ poṭṭivīṇu.} — In the wind a branch snapped and fell.
  \item \emph{Recall:} say the Malayalam for a heatwave, then the Malayalam for a cyclone, then say \emph{poṭṭivīḻuka} again
\end{itemize}

\begin{tcolorbox}[breakable,colback=teal!4,colframe=teal!35!black,title={Before you move on}]
What does \ml{പൊട്ടിവീഴുക} mean? (\textquotedblleft{}To snap and fall\textquotedblright{}.) Say it once more.
\end{tcolorbox}
